\section{Proof of Theorem \ref{theoremb}}\label{prooftheoremb}
The inequality \eqref{LyapleqFalc},
\(\sup_{\mu \in \mathcal{M}}\dim_{LY}^P(\rho_\ast \mu) \le \dim_F^P(\rho)\),
was proved in Section \ref{proofLyap}.

\subsection{Proof of \(\dim_F^P(\rho)\le \sup_{\mu \in \mathcal{M}}\dim_{LY}^P(\rho_\ast \mu) \)}
We recall that the limit cone of \(\rho\) is defined as
\[
\mathcal{L}_{\rho} = \lim_{T \to +\infty}\frac{1}{T} (\log S (\rho (\Gamma)),
\]
in the Hausdorff topology. It is a closed convex cone in \(\aaa^+\).
By the \(P\)-Anosov property if \(a \in \mathcal{L}_{\rho} \setminus \lbrace 0\rbrace\) we obtain
\[
\alpha_{i,j}(a) > 0\text{ for all }(i,j) \in S(P).
\]
Assume that \(r < r' < \dim_F^P(\rho)\) so that the series \eqref{falconerseries} diverges at \(r'\). By compactness there exists \(a \in \mathcal{L}_{\rho}\setminus \lbrace 0\rbrace\) such that for all open cones \(\mathcal{C}\) containing \(a\) we have
\[
\sum_{\g \in \G: \log S (\rho(\gamma)) \in \mathcal{C}} \exp \bigl(-F_{r'}^P\circ \log S (\rho(\gamma)) \bigr) = +\infty.
\]
We fix such a choice of \(a\) and notice that \(\psi(a) \ge F_{r'}^P(a)\) by \cite[Lem.\,III.1.3]{Qui02a} and \cite[Lem.\,4.2]{sambarino}. Thus, by Lemma \ref{lyapunovvsfalconer}, \(\psi(a) > F_{r}^P(a)\). We~will prove:
\begin{proposition}\label{goodmeasureslemma}
There exists a sequence \(T_k \to +\infty\) and \(\mu_k \in \M\) such that \(h_{\mu_k} = T_k \psi(a) + o(T_k)\) and \(\lambda(\rho_*\mu_k) = T_k a + o(T_k)\) when \(k \to +\infty\).
\end{proposition}

Since \(\psi(a) > F_r^P(a)\), we have from Lemma \ref{lyapunovvsfalconer} that \(L_{\psi(a)}^P(a) > r\).
Assuming Proposition \ref{goodmeasureslemma}, we obtain
\begin{align*}\dim_{LY}^P(\rho_*\mu_k) &= L_{h_{\mu_k}}^P(\rho_*\mu_k) = L_{T_k\psi(a) + o(T_k)}^P(T_k a + o(T_k))
\\ &= L_{\psi(a) + o(1)}^P(a + o(1)) = L_{\psi(a)}^P(a) + o(1) > r,\end{align*}
for \(k\) large enough. Which concludes the proof of the inequality and of Theorem~\ref{theoremb}.\qed

\subsection{Proof of Proposition \ref{goodmeasureslemma}}

Fix \(\beta:\aaa \to \R\) linear such that \(\beta(a) = F_r^P(a)\), and a decreasing sequence of open cones \(\mathcal{C}_k\), \(k = 1,2,3,\ldots\) whose intersection is \(\R_+ a\).

For a given \(k\) set\vspace*{-3pt}
\[
\psi_k(a) := \limsup_{T \to +\infty}\frac{1}{T}\log \#\bigl\lbrace \gamma \in \Gamma: \log S(\rho(\gamma)) \in \mathcal{C}_k\cap \lbrace \beta \le T \beta(a)\rbrace \bigr\rbrace. 
\]
We have \(\psi (a) = \lim_{k \to \infty } \psi_k (a) \) and\vspace*{-3pt}
\[
\psi_k(a) = \limsup_{T \to +\infty}\frac{1}{T}\log \#\bigl\lbrace \gamma \in \Gamma: \log S(\rho(\gamma)) \in \mathcal{C}_k \cap \lbrace (T-1)\beta(a) \le \beta \le T \beta(a)\rbrace \bigr\rbrace.
\]
Hence, we~may choose \(T_k \to +\infty\) such that the sets\vspace*{-3pt}
\[
A_k = \bigl\lbrace \gamma \in \G: \log S (\rho(\gamma))\in \mathcal{C}_k \cap \lbrace (T_k-1)\beta(a) \le \beta \le T_k \beta(a)\rbrace\bigr\rbrace,
\]
satisfy\vspace*{-3pt}
\[
\psi(a) = \lim_{k \to +\infty}\frac{1}{T_k}\log \# A_k.
\]
Moreover, by the Anosov property, there are constants \(C_1,C_2\) such that for \(\g \in A_k\), \(C_1T_k \le |\g| \le C_2T_k\). So there exist \(\ell_k \to \infty \) as \(k \to \infty \), such that the sets
\[
A'_k = \bigl\lbrace \gamma \in \G: \log S (\rho(\gamma)) \in \mathcal{C}_k \cap \lbrace (T_k-1)\beta(a) \le \beta \le T_k \beta(a)\rbrace,\, |\g| = \ell_k\bigr\rbrace,
\]
still satisfy\vspace*{-3pt}
\[
\psi(a) = \lim_{k \to +\infty}\frac{1}{T_k}\log \# A'_k.
\]

We now fix a semi-group \(S\) of \(\Gamma\), a finite subset \(F\), and maps \(L,R:\Gamma \to F\) given by Proposition \ref{semigrouplemma}. Observe that \(\rho(S)\) is Zariski dense since the Zariski closure of a semi-group is a group(see \cite{goldsheid-margulis}) and since \(S\) generates \(\Gamma\).

Let \(B_k \subset S\) be defined by\vspace*{-3pt}
\[
B_k = \lbrace L(\gamma)\gamma R(\gamma): \gamma \in A'_k\rbrace.
\]
Since \(F\) is finite, there exists \(C > 0\) such that\vspace*{-3pt}
\begin{align*} \# B_k &\ge C^{-1} \# A'_k, & | |\g| - \ell_k | &\le C {\textrm { for }} \g \in B_k \end{align*}
and fixing some norm on \(\aaa\) we have\vspace*{-3pt}
\[
\|\log S(\rho(\gamma)) - \log S(\rho(L(\gamma)\gamma R(\gamma)))\| \le C,
\]
for all \(\gamma \in \Gamma\).

Fix a parameter \(\e> 0 \), we~say a splitting \(\R^d = E_1 \oplus \cdots \oplus E_d\) into one dimensional subspaces is \(\e\)-non-degenerate if
the angle between \(E_i\) and \(\bigoplus_{j \neq i}E_j\) is at least~\(\e\) for all~\(i\). We~say two splittings \((E_1,\ldots, E_d), (E_1',\ldots,E_d')\) are \(\e\)-close if the angle between~\(E_i\) and \(E_i'\) is at most \(\e\) for all \(i\).

We say that an element \(g \in \SL(d,\R)\) is \(\e\)-diagonalizable if it has \(d\) distinct real eigenvalues \(|\lambda_1| > \cdots > |\lambda_d|\) satisfying \(|\lambda_i|/|\lambda_{i+1}| > e^{\e}\) for \(i = 1,\ldots, d-1\), and the corresponding splitting into eigenspaces is \(\e\)-non-degenerate.

Since \(\rho(S)\) is Zariski dense we obtain from \cite[Th.\,6.8]{AMS95}:
\begin{lemma}[Abels-Margulis-Soifer]\label{abslemma}
There exists \(\e_0 > 0\) and a finite set \(S_0 \subset S\) such that for all \(\gamma \in \Gamma\) the set \(\rho(\gamma S_0)\) contains at least one \(\e_0\)-diagonalizable element.
\end{lemma}

For what follows we fix \(S_0\) and \(\e_0 > 0\) given by the previous lemma. We~also fix \(\delta_0 > 0\) given by the following:
\begin{lemma}
There exists \(\delta_0 > 0\) such that if \(B \) is a set of \(\e_0\)-diagonalizable elements of \(\SL(d,\R) \) whose eigenspace splittings are \(\delta_0\)-close, then
\(\gamma_1\cdots \gamma_n\) is \(\e_0/2\)-diagonalizable for all \(n\) and \(\gamma_1,\ldots, \gamma_n \in B\).
\end{lemma}
\begin{proof}
This follows directly from \cite[Lem.\,2.17]{2021breuillard-sert}.
\end{proof}

We now refine the family \(B_k\) so as to obtain images under \(\rho\) which are diagonalizable with nearby splittings.
\begin{lemma}\label{bklemma}
There exists a decreasing positive sequence \(\epsilon_k \downarrow 0\) and a sequence \(\ell '_k \uparrow \infty \) such that for all \(k\) large there exists \(C_k \subset B_k\cdot S_0\) with the following properties:
\begin{enumerate}
\item\label{bklemma1} \(C_k\) freely generates a free semi-group in \(\Gamma\).
\item\label{bklemma2} For all \(\gamma \in C_k\) the element \(\rho(\gamma)\) is \(\e_0\)-diagonalizable.
\item\label{bklemma3} For all \(\gamma_1,\gamma_2 \in C_k\) the eigenspace splittings of \(\rho(\gamma_1)\) and \(\rho(\gamma_2)\) are \(\delta_0\)-close.
\item\label{bklemma4} One has \(\psi(a) = \lim_{k\to +\infty}\frac{1}{T_k}\log \# C_k\).
\item\label{bklemma5} \(|\g| = \ell '_k \) for \(\g \in C_k \).
\item\label{bklemma6} For all \(\gamma \in C_k\) one has \(\|(\log S(\rho(\gamma)) - T_k a\| \le \epsilon_k T_k\).
\end{enumerate}
\end{lemma}
\begin{proof}
The space of \(\e_0\)-non-degenerate splittings is compact and hence we may cover it by some finite number \(N\) of subsets with the property that if two splittings belong to the same subset they are \(\delta_0\)-close.

From Lemma \ref{abslemma} it follows that for all \(\gamma \in B_k\) there exists \(f_\gamma \in S_0\) such that \(\rho(\gamma f_\gamma)\) is \(\e_0\)-diagonalizable. By the pigeonhole principle we may choose a set \(C'_k\) of at least \(\# B_k /N\) of the elements \(\gamma f_{\gamma}\) such that the eigenspace splitting of \(\rho\) applied to any two are \(\delta_0\)-close.
Hence, we~have constructed \(C'_k\) with properties \eqref{bklemma2}, \eqref{bklemma3} and \eqref{bklemma4} above.

Since \(F \cup S_0\) is finite there exists \(C > 0\) such that, for all \(k\) and \(\gamma \in A_k\),
\begin{eqnarray*} | |L(\gamma)\gamma R(\gamma) f_{L(\gamma)\gamma R(\gamma)}| - \ell_k | &\le &C \; {\textrm { and }}\\
\|(\log S(\rho(L(\gamma)\gamma R(\gamma) f_{L(\gamma)\gamma R(\gamma)})) - (\log S(\rho(\gamma))\| &\le &C.\end{eqnarray*}

By the pigeonhole principle again, we~can find \(\ell'_k, | \ell '_k - \ell_k | \le C\), such that \(C''_k : = \g \in C'_k, |\g | = \ell '_k \) satisfies properties \eqref{bklemma2}, \eqref{bklemma3}, \eqref{bklemma4} and \eqref{bklemma5}. Moreover,
if \(\gamma \in A_k\) then by definition we have
\[
\Bigl(\frac{1}{T_k}(\log S(\rho(\gamma))- a\Bigr) \in \mathcal{C}_k \cap \Bigl\lbrace \frac{T_k - 1}{T_k}\beta(a) \le \beta \le \beta(a)\Bigr\rbrace,
\]
hence setting \(\epsilon_k\) to be \(C/T_k\) plus the diameter of the set on the right-hand side, we~have
\[
\|(\log S(\rho(L(\gamma)\gamma R(\gamma) f_\gamma))- T_k a\| \le \epsilon_k T_k,
\]
for all \(\gamma \in A_k\), which establishes property (6) for \(C'_k\) and \(C''_k\).

To establish \eqref{bklemma1} we notice that \(S\) is a quasi-geodesic semi-group. Therefore, by~the Morse lemma, there exists some \(R > 0\) such that if \(\gamma_1,\ldots, \gamma_n, \eta_1,\ldots,\eta_n \in S\) and \(\gamma_1\cdots \gamma_n = \eta_1\cdots \eta_n\), then there is a geodesic ray \(\alpha = \{\alpha_n \}_{ n \ge 0}\), such that for all \(j \in [0, n]\), \begin{equation}\label{tracking}\dist(\gamma_1\cdots \gamma_j, \alpha), \, \dist (\eta_1\cdots \eta_j, \alpha) < R.\end{equation}

We claim that if we assume that the elements of \(C_k \subset C''_k \) are \(6R\)-separated, equation \eqref{tracking} forces \(\g_j = \eta_j \) for all \(j \in [0, n]\).
Hence it suffices to refine \(C''_k\) so that it is a \(6R\)-separated set to obtain that it freely generates a free sub-semi-group of \(\Gamma\). Doing this decreases the number of elements at most by a factor equal to the number of elements in a ball of radius \(6R\) in \(\Gamma\), so property \eqref{bklemma4} still holds.

To prove the claim, we~prove by induction on \(j\) that \(\g_i = \eta_i \) for \(i \le j\). We~have \(\g_0 = \eta_0 = e\). Assume that \(\g_i = \eta_i \) for \(i \le j\). Then
\(\gamma_1\cdots \gamma_j = \eta_1\cdots \eta_j\) and there exist \(N_1, N_2\) and \(N_3\) such that
\[
 \dist (\gamma_1\cdots \gamma_j, \alpha_{N_1}),\;\dist (\gamma_1\cdots \gamma_j\g_{j+1}, \alpha_{N_2}) {\textrm { and }}\dist (\gamma_1\cdots \gamma_j\eta_{j+1}, \alpha_{N_3}) \le R.
\]
Setting \(\beta_t := (\gamma_1\cdots \gamma_j)^{-1} \alpha_{N_t} \) for \(t = 1,2,3\), we have \(\beta_{N_1}, \beta_{N_2}\) and \(\beta_{N_3} \) aligned on the same geodesic and satisfying:
\[
 |\beta_1|,\, \dist (\g_{j+1}, \beta_2),\, \dist (\eta_{j+1}, \beta_3) < R \quand |\g_{j+1}|= |\eta_{j+1}| = \ell'_k.
\]
If \(\ell'_k\) is large enough, \(\beta_1 \) is outside the interval \([\beta_2, \beta_3]\) (see Lemma \ref{boundeddistanelemma}) and
\[
 \dist (\beta_1, \beta_2),\, \dist (\beta_1, \beta_3) \; \in \; [\ell'_k -2R, \ell'_k +2R ].
\]
It follows that \(\dist (\beta_2, \beta_3) \le 4R\) and therefore that \(\dist (\alpha_{j+1}, \eta_{j+1}) < 6R\). Since \(C_k\) is \(6R\)-separated, \(\alpha_{j+1} = \eta_{j+1}\) as claimed.
\end{proof}

\begin{corollary}
If \(\mu_k\) is the uniform probability measure on \(C_k\) then one has
\[
\lim_{k \to +\infty}\frac{1}{T_k}h_{\mu_k} = \psi(a).
\]
\end{corollary}
\begin{proof}
This follows immediately from properties \eqref{bklemma1} and \eqref{bklemma4} of the previous lemma.
\end{proof}

We now use the fact that for \(\e_0\)-diagonalizable elements with close splittings the singular values and eigenvalues are close and they almost multiply under composition.
\begin{lemma}
There exists a decreasing positive sequence \(\epsilon_k' \downarrow 0\) such that for all \(k\) large enough, all \(n\) and \(\gamma_1,\ldots, \gamma_n \in C_k\) one has
\[
\Bigl\|\frac{1}{n}(\log S(\rho(\gamma_1\cdots \gamma_n)) - T_k a\Bigr\| \le \epsilon_k' T_k.
\]
\end{lemma}

\begin{proof}
Let \(L(g)\) denote the vector of absolute values of eigenvalues of a diagonalizable element \(g \in \SL(d,\R)\) in decreasing order.
By \cite[Lem.\,2.16]{2021breuillard-sert} there exists \(C > 0\) such that for all \(\e_0/2\)-diagonalizable \(g \in \SL(d,\R)\) one has
\[
\|\log S(g) - \log L(g)\| \le C.
\]
By \cite[Prop.\,2.17]{2021breuillard-sert} there exists \(C > 0\) independent of \(k\) and \(n\) such that we have
\[
\|\ts (\log L(\rho(\gamma_1\cdots \gamma_n)) - \sum_{i = 1}^n (\log L (\rho(\gamma_i))\| \le Cn.
\]
Combining this with Lemma \ref{bklemma} we obtain
\begin{align*}\|(\log S(\rho(\gamma_1 \cdots \gamma_n)) - n T_k a\| &\le C + \|(\log L (\rho(\gamma_1\cdots \gamma_n))- nT_k a\|
\\ &\ts\le C(n+1) + \|\sum_{i = 1}^n(\log L (\rho(\gamma_i)) - nT_k a\|
\\ &\ts\le C(2n+1) + \|\sum_{i = 1}^n(\log S (\rho(\gamma_i))- nT_k a\|
\\ &\le C(2n+1) + n \epsilon_k T_k.
\end{align*}
This proves the desired result with \(\epsilon_k' = 3C / T_k + \epsilon_k\).
\end{proof}

The following immediate corollary concludes the proof of Proposition \ref{goodmeasureslemma}.

\begin{corollary}
For any sequence of probability measures \(\mu_k\) with \(\mu_k(C_k) = 1\) one has
\[
\lim_{k \to +\infty}\frac{1}{T_k}\lambda(\rho_*\mu_k) = a.
\]
\end{corollary}

